\textbf{COLUNI 2025 -- Questão 09}\par\medskip Leia o excerto a seguir.

\par

No começo, andavam pra cima e pra baixo, o homem tinha os olhos grandes para as coisas que via na fazenda, e a mulher conseguia fingir interesse. Era \underline{atabalhoada}, entrava nos lugares sem saber se poderia entrar, repetia as mesmas frases de espanto, sorria de forma quase discreta do que chamava de ignorância do povo, quando se dispunha a perguntar e obtinha uma resposta diferente da que presumia ser verdadeira. (VIEIRA JUNIOR, Itamar. \textit{Torto arado}. São Paulo: Todavia, 2019. p. 210)

\par

Marque a alternativa que apresenta a palavra que substitui, ADEQUADAMENTE, o termo grifado em seu efeito de sentido no contexto da obra.

\par

\begin{enumerate}
\item alçada.
\item incomodada.
\item desordenada.
\item irritada.
\end{enumerate}
